% !TeX encoding = UTF-8
w\documentclass[aspectratio=169]{beamer}

\usetheme{metropolis}
\usepackage[spanish]{babel}
\usepackage[utf8]{inputenc}
\usepackage[T1]{fontenc}
\usepackage{tikz}
\usepackage{booktabs}
\usepackage{amsmath}
\usepackage{hyperref}

\title{Sistemas Distribuidos}
\subtitle{Clase 2: Modelos, fallas, tiempo y comunicación}
\author{Rodolfo L. Sumoza}
\institute{Universidad Austral}
\date{}

\begin{document}

% ------------------------------------------------
\begin{frame}
  \titlepage
\end{frame}

% ------------------------------------------------
\begin{frame}{La pregunta central de hoy}

\Large

\begin{center}
\textbf{¿Cómo diseñamos un sistema confiable cuando no podemos confiar en la red, en el tiempo ni en los nodos?}
\end{center}

\vspace{1cm}

\normalsize

\begin{itemize}
    \item Los mensajes pueden perderse.
    \item Los nodos pueden fallar.
    \item La red puede particionarse.
    \item No existe un reloj global perfecto.
\end{itemize}

\end{frame}

% ------------------------------------------------
\begin{frame}{Un sistema distribuido no falla como una computadora}

\Large

\begin{center}
\textbf{En una máquina local, el fallo suele ser visible.}
\end{center}

\vspace{0.5cm}

\begin{center}
\textbf{En un sistema distribuido, el fallo puede parecer demora.}
\end{center}

\vspace{0.8cm}

\normalsize

\begin{block}{Problema fundamental}
Si un nodo no responde, no sabemos si:
\end{block}

\begin{itemize}
    \item falló;
    \item está lento;
    \item el mensaje se perdió;
    \item la red está particionada;
    \item la respuesta todavía no llegó.
\end{itemize}

\end{frame}

% ------------------------------------------------
\begin{frame}{Modelo mental de un sistema distribuido}

\begin{center}
\begin{tikzpicture}[
    node/.style={circle, draw, thick, minimum size=1.3cm},
    msg/.style={->, thick}
]

\node[node] (A) at (0,0) {A};
\node[node] (B) at (4,1.5) {B};
\node[node] (C) at (4,-1.5) {C};
\node[node] (D) at (8,0) {D};

\draw[msg] (A) -- (B);
\draw[msg] (A) -- (C);
\draw[msg] (B) -- (D);
\draw[msg] (C) -- (D);
\draw[msg, dashed] (B) -- (C);

\end{tikzpicture}
\end{center}

\vspace{0.5cm}

\Large
\begin{center}
\textbf{El sistema es una colección de procesos que se comunican mediante mensajes.}
\end{center}

\end{frame}

% ------------------------------------------------
\begin{frame}{Tres dimensiones del problema}

\Large

\begin{enumerate}
    \item \textbf{Comunicación}
    
    ¿Los mensajes llegan? ¿En qué orden? ¿Con qué demora?
    
    \vspace{0.4cm}
    
    \item \textbf{Tiempo}
    
    ¿Podemos medir el tiempo real? ¿Los relojes coinciden?
    
    \vspace{0.4cm}
    
    \item \textbf{Fallas}
    
    ¿Qué tipo de fallas esperamos tolerar?
\end{enumerate}

\end{frame}

% ------------------------------------------------
\begin{frame}{Modelos de comunicación}

\begin{table}
\centering
\begin{tabular}{lll}
\toprule
\textbf{Modelo} & \textbf{Supuesto} & \textbf{Dificultad} \\
\midrule
Síncrono & Tiempo acotado & Más fácil \\
Asíncrono & Sin límite temporal & Más realista \\
Parcialmente síncrono & Eventualmente estable & Muy usado \\
\bottomrule
\end{tabular}
\end{table}

\vspace{0.7cm}

\begin{block}{Idea clave}
El modelo elegido determina qué algoritmos son posibles.
\end{block}

\end{frame}

% ------------------------------------------------
\begin{frame}{Sistema síncrono}

\Large

\begin{center}
\textbf{Existe un límite conocido para la demora de mensajes y procesamiento.}
\end{center}

\vspace{0.7cm}

\normalsize

\begin{itemize}
    \item Podemos usar timeouts confiables.
    \item Podemos distinguir demora de fallo.
    \item Es un modelo más simple.
    \item Pero suele ser poco realista en Internet.
\end{itemize}

\end{frame}

% ------------------------------------------------
\begin{frame}{Sistema asíncrono}

\Large

\begin{center}
\textbf{No hay límite conocido para la demora.}
\end{center}

\vspace{0.7cm}

\normalsize

\begin{itemize}
    \item Un mensaje puede tardar arbitrariamente.
    \item Un nodo lento puede parecer muerto.
    \item Los timeouts son sospechas, no certezas.
    \item Muchos problemas se vuelven imposibles bajo fallas.
\end{itemize}

\end{frame}

% ------------------------------------------------
\begin{frame}{El problema del timeout}

\Large

\begin{center}
\textbf{Si B no responde, ¿B falló o la red está lenta?}
\end{center}

\vspace{0.6cm}

\begin{center}
\begin{tikzpicture}[
    node/.style={circle, draw, thick, minimum size=1.3cm},
    msg/.style={->, thick}
]

\node[node] (A) at (0,0) {A};
\node[node] (B) at (6,0) {B};

\draw[msg] (A) -- node[above]{request} (B);
\draw[msg, dashed] (B) -- node[below]{response?} (A);

\end{tikzpicture}
\end{center}

\vspace{0.5cm}

\normalsize

\begin{block}{Conclusión}
En sistemas distribuidos, un timeout no prueba una falla. Solo expresa una sospecha.
\end{block}

\end{frame}

% ------------------------------------------------
\begin{frame}{Tipos de fallas}

\begin{table}
\centering
\begin{tabular}{lll}
\toprule
\textbf{Tipo de falla} & \textbf{Descripción} & \textbf{Ejemplo} \\
\midrule
Crash & El nodo se detiene & Servidor caído \\
Omisión & No envía o no recibe & Paquete perdido \\
Timing & Responde tarde & Timeout \\
Bizantina & Comportamiento arbitrario & Nodo malicioso \\
\bottomrule
\end{tabular}
\end{table}

\vspace{0.5cm}

\begin{block}{Regla práctica}
Mientras más fuerte es el modelo de falla, más costoso es el algoritmo.
\end{block}

\end{frame}

% ------------------------------------------------
\begin{frame}{Falla crash}

\Large

\begin{center}
\textbf{El nodo simplemente deja de funcionar.}
\end{center}

\vspace{0.7cm}

\normalsize

\begin{itemize}
    \item Es el modelo de falla más simple.
    \item El nodo no envía información incorrecta.
    \item Muchos protocolos prácticos asumen este modelo.
\end{itemize}

\vspace{0.5cm}

\begin{exampleblock}{Ejemplo}
Un servidor de base de datos se apaga inesperadamente.
\end{exampleblock}

\end{frame}

% ------------------------------------------------
\begin{frame}{Falla bizantina}

\Large

\begin{center}
\textbf{El nodo puede comportarse de manera arbitraria.}
\end{center}

\vspace{0.7cm}

\normalsize

\begin{itemize}
    \item Puede mentir.
    \item Puede enviar mensajes distintos a distintos nodos.
    \item Puede actuar de forma maliciosa o corrupta.
    \item Requiere protocolos mucho más costosos.
\end{itemize}

\end{frame}

% ------------------------------------------------
\begin{frame}{El tiempo como problema distribuido}

\Large

\begin{center}
\textbf{No existe un reloj global perfecto.}
\end{center}

\vspace{0.7cm}

\normalsize

\begin{itemize}
    \item Cada máquina tiene su propio reloj.
    \item Los relojes derivan.
    \item La sincronización nunca es perfecta.
    \item Ordenar eventos es más difícil de lo que parece.
\end{itemize}

\end{frame}

% ------------------------------------------------
\begin{frame}{Orden físico vs orden causal}

\Large

\begin{center}
\textbf{No siempre necesitamos saber qué ocurrió primero en tiempo real.}
\end{center}

\vspace{0.6cm}

\begin{center}
\textbf{Muchas veces necesitamos saber qué pudo haber causado qué.}
\end{center}

\vspace{0.8cm}

\normalsize

\begin{block}{Relación happens-before}
Un evento A ocurre antes que B si A pudo haber influido causalmente en B.
\end{block}

\end{frame}

% ------------------------------------------------
\begin{frame}{Ejemplo de causalidad}

\begin{center}
\begin{tikzpicture}[
    event/.style={circle, fill=black, minimum size=5pt, inner sep=0pt},
    msg/.style={->, thick},
    proc/.style={thick}
]

\draw[proc] (0,2) -- (9,2);
\draw[proc] (0,0) -- (9,0);

\node[left] at (0,2) {P1};
\node[left] at (0,0) {P2};

\node[event] (a) at (2,2) {};
\node[event] (b) at (5,2) {};
\node[event] (c) at (7,0) {};

\draw[msg] (b) -- (c);

\node[above] at (2,2) {a};
\node[above] at (5,2) {b};
\node[below] at (7,0) {c};

\end{tikzpicture}
\end{center}

\vspace{0.5cm}

\Large

\begin{center}
\textbf{Si b envía un mensaje que recibe c, entonces b ocurrió antes que c.}
\end{center}

\end{frame}

% ------------------------------------------------
\begin{frame}{Relojes lógicos}

\Large

\begin{center}
\textbf{Un reloj lógico no mide tiempo físico.}
\end{center}

\vspace{0.6cm}

\begin{center}
\textbf{Mide orden causal.}
\end{center}

\vspace{0.8cm}

\normalsize

\begin{itemize}
    \item Cada proceso mantiene un contador.
    \item El contador aumenta con cada evento.
    \item Los mensajes transportan marcas temporales.
    \item Al recibir un mensaje, el proceso actualiza su contador.
\end{itemize}

\end{frame}

% ------------------------------------------------
\begin{frame}{Reglas de Lamport}

\Large

\begin{enumerate}
    \item Antes de cada evento local:
    \[
    C_i = C_i + 1
    \]
    
    \item Al enviar un mensaje:
    \[
    send(m, C_i)
    \]
    
    \item Al recibir un mensaje con timestamp $t$:
    \[
    C_i = \max(C_i, t) + 1
    \]
\end{enumerate}

\end{frame}

% ------------------------------------------------
\begin{frame}{Qué garantizan los relojes de Lamport}

\Large

\begin{block}{Garantía}
Si un evento $a$ ocurrió antes que un evento $b$, entonces:
\[
C(a) < C(b)
\]
\end{block}

\vspace{0.7cm}

\begin{alertblock}{Pero cuidado}
Si:
\[
C(a) < C(b)
\]
no necesariamente significa que $a$ causó a $b$.
\end{alertblock}

\end{frame}

% ------------------------------------------------
\begin{frame}{Trade-off permanente}

\Large

\begin{center}
\textbf{Todo sistema distribuido vive negociando entre tres fuerzas:}
\end{center}

\vspace{0.7cm}

\begin{itemize}
    \item rendimiento;
    \item consistencia;
    \item tolerancia a fallas.
\end{itemize}

\vspace{0.7cm}

\begin{center}
\textbf{No se puede maximizar todo al mismo tiempo.}
\end{center}

\end{frame}

% ------------------------------------------------
\begin{frame}{Actividad en clase}

\Large

\begin{center}
\textbf{Diseñemos un sistema de pagos distribuido.}
\end{center}

\vspace{0.7cm}

\normalsize

Preguntas:

\begin{enumerate}
    \item ¿Qué pasa si se cae un datacenter?
    \item ¿Qué pasa si dos regiones ven saldos distintos?
    \item ¿Preferimos rechazar pagos o aceptar inconsistencia temporal?
    \item ¿Qué modelo de falla asumimos?
    \item ¿Qué garantías prometemos al usuario?
\end{enumerate}

\end{frame}

% ------------------------------------------------
\begin{frame}{Lo que tienen que recordar}

\Large

\begin{enumerate}
    \item Un sistema distribuido falla parcialmente.
    \item Un timeout no prueba una falla.
    \item El modelo de comunicación cambia lo que es posible.
    \item No hay reloj global perfecto.
    \item El orden causal importa más que el tiempo físico.
    \item Toda solución implica trade-offs.
\end{enumerate}

\end{frame}

% ------------------------------------------------
\begin{frame}{Cierre}

\Large

\begin{center}
\textbf{La pregunta no es si el sistema va a fallar.}
\end{center}

\vspace{0.8cm}

\begin{center}
\textbf{La pregunta es qué garantías seguirá ofreciendo cuando falle.}
\end{center}

\end{frame}

\end{document}
